
\documentclass[11pt,a4paper]{article}
% Shared preamble for Wolf lecture note transcriptions
% Usage: % Shared preamble for Wolf lecture note transcriptions
% Usage: % Shared preamble for Wolf lecture note transcriptions
% Usage: \input{preamble} in each chapter file
% Include guard: safe to \input multiple times (e.g. from main_wolf.tex)
\ifdefined\wolfpreambleloaded\endinput\fi
\def\wolfpreambleloaded{}

\usepackage{amsmath,amssymb,amsthm}
% Match Wolf book-class layout: a4paper, ~345pt text width
\usepackage[a4paper, textwidth=345pt, textheight=550pt]{geometry}
\usepackage{hyperref}
\usepackage{mathtools}

% ── Notation (consistent across all chapters) ──
\newcommand{\mcM}{\mathcal{M}}       % matrix algebra M_d(C)
\newcommand{\mcB}{\mathcal{B}}       % bounded operators B(H)
\newcommand{\mcA}{\mathcal{A}}       % multiplicative domain
\newcommand{\mcH}{\mathcal{H}}       % Hilbert space H
\newcommand{\diag}{\operatorname{diag}}
\newcommand{\one}{\mathbb{1}}        % identity operator
\newcommand{\CC}{\mathbb{C}}
\newcommand{\RR}{\mathbb{R}}
\newcommand{\NN}{\mathbb{N}}
\newcommand{\ZZ}{\mathbb{Z}}
\newcommand{\tr}{\operatorname{tr}}
\newcommand{\spec}{\operatorname{spec}}
\newcommand{\rank}{\operatorname{rank}}
\newcommand{\supp}{\operatorname{supp}}
\newcommand{\range}{\operatorname{range}}
\newcommand{\spn}{\operatorname{span}}
\newcommand{\FT}{\mathcal{F}_T}      % fixed point set
\newcommand{\XT}{\mathcal{X}_T}      % peripheral eigenspace
\newcommand{\mcX}{\mathcal{X}}       % peripheral eigenspace (bare)
\newcommand{\id}{\mathrm{id}}        % identity map
\newcommand{\ket}[1]{|#1\rangle}
\newcommand{\bra}[1]{\langle#1|}
\newcommand{\braket}[2]{\langle#1|#2\rangle}
\newcommand{\ketbra}[2]{|#1\rangle\!\langle#2|}

% ── Helper: properly undefine a theorem environment for redefinition ──
% Usage: \UndefEnv{theorem} before \newtheorem{theorem}{...}
\makeatletter
\newcommand{\UndefEnv}[1]{%
  % Remove from section reset list (if registered there)
  \@removefromreset{#1}{section}%
  \expandafter\let\csname #1\endcsname\relax
  \expandafter\let\csname end#1\endcsname\relax
  \expandafter\let\csname c@#1\endcsname\relax
}
\makeatother

% ── Theorem environments ──
\theoremstyle{plain}
\newtheorem{theorem}{Theorem}[section]
\newtheorem{proposition}[theorem]{Proposition}
\newtheorem{corollary}[theorem]{Corollary}
\newtheorem{lemma}[theorem]{Lemma}
\theoremstyle{definition}
\newtheorem{definition}[theorem]{Definition}
\newtheorem{example}[theorem]{Example}
\newtheorem{problem}[theorem]{Problem}
\theoremstyle{remark}
\newtheorem{remark}[theorem]{Remark} in each chapter file
% Include guard: safe to \input multiple times (e.g. from main_wolf.tex)
\ifdefined\wolfpreambleloaded\endinput\fi
\def\wolfpreambleloaded{}

\usepackage{amsmath,amssymb,amsthm}
% Match Wolf book-class layout: a4paper, ~345pt text width
\usepackage[a4paper, textwidth=345pt, textheight=550pt]{geometry}
\usepackage{hyperref}
\usepackage{mathtools}

% ── Notation (consistent across all chapters) ──
\newcommand{\mcM}{\mathcal{M}}       % matrix algebra M_d(C)
\newcommand{\mcB}{\mathcal{B}}       % bounded operators B(H)
\newcommand{\mcA}{\mathcal{A}}       % multiplicative domain
\newcommand{\mcH}{\mathcal{H}}       % Hilbert space H
\newcommand{\diag}{\operatorname{diag}}
\newcommand{\one}{\mathbb{1}}        % identity operator
\newcommand{\CC}{\mathbb{C}}
\newcommand{\RR}{\mathbb{R}}
\newcommand{\NN}{\mathbb{N}}
\newcommand{\ZZ}{\mathbb{Z}}
\newcommand{\tr}{\operatorname{tr}}
\newcommand{\spec}{\operatorname{spec}}
\newcommand{\rank}{\operatorname{rank}}
\newcommand{\supp}{\operatorname{supp}}
\newcommand{\range}{\operatorname{range}}
\newcommand{\spn}{\operatorname{span}}
\newcommand{\FT}{\mathcal{F}_T}      % fixed point set
\newcommand{\XT}{\mathcal{X}_T}      % peripheral eigenspace
\newcommand{\mcX}{\mathcal{X}}       % peripheral eigenspace (bare)
\newcommand{\id}{\mathrm{id}}        % identity map
\newcommand{\ket}[1]{|#1\rangle}
\newcommand{\bra}[1]{\langle#1|}
\newcommand{\braket}[2]{\langle#1|#2\rangle}
\newcommand{\ketbra}[2]{|#1\rangle\!\langle#2|}

% ── Helper: properly undefine a theorem environment for redefinition ──
% Usage: \UndefEnv{theorem} before \newtheorem{theorem}{...}
\makeatletter
\newcommand{\UndefEnv}[1]{%
  % Remove from section reset list (if registered there)
  \@removefromreset{#1}{section}%
  \expandafter\let\csname #1\endcsname\relax
  \expandafter\let\csname end#1\endcsname\relax
  \expandafter\let\csname c@#1\endcsname\relax
}
\makeatother

% ── Theorem environments ──
\theoremstyle{plain}
\newtheorem{theorem}{Theorem}[section]
\newtheorem{proposition}[theorem]{Proposition}
\newtheorem{corollary}[theorem]{Corollary}
\newtheorem{lemma}[theorem]{Lemma}
\theoremstyle{definition}
\newtheorem{definition}[theorem]{Definition}
\newtheorem{example}[theorem]{Example}
\newtheorem{problem}[theorem]{Problem}
\theoremstyle{remark}
\newtheorem{remark}[theorem]{Remark} in each chapter file
% Include guard: safe to \input multiple times (e.g. from main_wolf.tex)
\ifdefined\wolfpreambleloaded\endinput\fi
\def\wolfpreambleloaded{}

\usepackage{amsmath,amssymb,amsthm}
% Match Wolf book-class layout: a4paper, ~345pt text width
\usepackage[a4paper, textwidth=345pt, textheight=550pt]{geometry}
\usepackage{hyperref}
\usepackage{mathtools}

% ── Notation (consistent across all chapters) ──
\newcommand{\mcM}{\mathcal{M}}       % matrix algebra M_d(C)
\newcommand{\mcB}{\mathcal{B}}       % bounded operators B(H)
\newcommand{\mcA}{\mathcal{A}}       % multiplicative domain
\newcommand{\mcH}{\mathcal{H}}       % Hilbert space H
\newcommand{\diag}{\operatorname{diag}}
\newcommand{\one}{\mathbb{1}}        % identity operator
\newcommand{\CC}{\mathbb{C}}
\newcommand{\RR}{\mathbb{R}}
\newcommand{\NN}{\mathbb{N}}
\newcommand{\ZZ}{\mathbb{Z}}
\newcommand{\tr}{\operatorname{tr}}
\newcommand{\spec}{\operatorname{spec}}
\newcommand{\rank}{\operatorname{rank}}
\newcommand{\supp}{\operatorname{supp}}
\newcommand{\range}{\operatorname{range}}
\newcommand{\spn}{\operatorname{span}}
\newcommand{\FT}{\mathcal{F}_T}      % fixed point set
\newcommand{\XT}{\mathcal{X}_T}      % peripheral eigenspace
\newcommand{\mcX}{\mathcal{X}}       % peripheral eigenspace (bare)
\newcommand{\id}{\mathrm{id}}        % identity map
\newcommand{\ket}[1]{|#1\rangle}
\newcommand{\bra}[1]{\langle#1|}
\newcommand{\braket}[2]{\langle#1|#2\rangle}
\newcommand{\ketbra}[2]{|#1\rangle\!\langle#2|}

% ── Helper: properly undefine a theorem environment for redefinition ──
% Usage: \UndefEnv{theorem} before \newtheorem{theorem}{...}
\makeatletter
\newcommand{\UndefEnv}[1]{%
  % Remove from section reset list (if registered there)
  \@removefromreset{#1}{section}%
  \expandafter\let\csname #1\endcsname\relax
  \expandafter\let\csname end#1\endcsname\relax
  \expandafter\let\csname c@#1\endcsname\relax
}
\makeatother

% ── Theorem environments ──
\theoremstyle{plain}
\newtheorem{theorem}{Theorem}[section]
\newtheorem{proposition}[theorem]{Proposition}
\newtheorem{corollary}[theorem]{Corollary}
\newtheorem{lemma}[theorem]{Lemma}
\theoremstyle{definition}
\newtheorem{definition}[theorem]{Definition}
\newtheorem{example}[theorem]{Example}
\newtheorem{problem}[theorem]{Problem}
\theoremstyle{remark}
\newtheorem{remark}[theorem]{Remark}

% --- Adjust numbering to match Wolf's Chapter-2 layout ---
\numberwithin{equation}{section}

% Wolf uses independent numbering for Theorems vs Propositions vs Examples,
% and Problems are globally numbered across the notes.
% The shared preamble defines these environments sharing one counter, so we
% redefine the ones needed here.
\UndefEnv{proposition}
\UndefEnv{example}
\UndefEnv{problem}

\theoremstyle{plain}
\newtheorem{proposition}{Proposition}[section]
\theoremstyle{definition}
\newtheorem{example}{Example}[section]
\newtheorem{problem}{Problem} % global numbering (no [section])

\title{Chapter 2: Representations \\ {\small (Wolf Lecture Notes - Transcription)}}
\date{}

\begin{document}
\maketitle
\tableofcontents

% Start chapter numbering at 2 (Chapter 2 as Section 2).
\setcounter{section}{1}

\section{Representations}

In this chapter we will discuss the most important ways of parameterizing and
representing a quantum channel:
(i) in terms of a bipartite quantum state, leading to the state-channel duality
introduced by Jamiolkowski and Choi,
(ii) as the reduced dynamics of a larger (unitarily evolving) system as expressed
by the representation theorems of Kraus, Stinespring, and Neumark (for POVMs),
and
(iii) as a linear map represented in terms of a ``transfer matrix''.
This will be supplemented with a brief discussion about normal forms.

\subsection{Jamiolkowski and Choi}

We saw in Prop.~1.2 that complete positivity of a linear map $T$ is equivalent to
positivity of the operator
\[
    \tau := (T \otimes \id)(\ketbra{\Omega}{\Omega})
\]
which is obtained by letting $T$ act on half of a maximally entangled state $\Omega$.
In fact, the operator $\tau$ obtained in this way encodes not only complete positivity
but every property of $T$. This reflects the simple equivalence
$\mcB(\mcM_d,\mcM_{d'}) \simeq \mcM_{dd'}$ which was certainly observed and
used at various places, but in the context of quantum channels the credit goes to
Choi and Jamiolkowski. $d\tau$ is often called \emph{Choi matrix}, and if $T$ is a
trace-preserving quantum channel, then $\tau$ is called the corresponding
\emph{Jamiolkowski state}.

\begin{proposition}[Choi--Jamiolkowski representation of maps]
    The following provides a one-to-one correspondence between linear maps
    $T \in \mcB(\mcM_d,\mcM_{d'})$ and operators $\tau \in \mcB(\CC^{d'}\otimes \CC^{d})$:
    \begin{equation}\label{eq:2.1}
        \begin{aligned}
            \tau           & = (T \otimes \id_d)(\ketbra{\Omega}{\Omega}), \\
            \tr\![A\,T(B)] & = d\,\tr\![\tau\,(A \otimes B^{T})],
        \end{aligned}
    \end{equation}
    for all $A \in \mcM_{d'}$, $B \in \mcM_d$ and $\Omega \in \CC^d\otimes \CC^d$
    being a maximally entangled state as in Eq.~(1.11). The maps $T \mapsto \tau$ and
    $\tau \mapsto T$ defined by \eqref{eq:2.1} are mutual inverses and lead to the
    following correspondences:
    \begin{itemize}
        \item Hermiticity: $\tau = \tau^\dagger$ iff $T(B^\dagger)=T(B)^\dagger$ for all
              $B \in \mcM_d$,\footnote{Note that this is in turn equivalent to $T$ being a
                  Hermiticity preserving map, i.e., $T(X)=T(X)^\dagger$ for all $X=X^\dagger$.}
        \item Complete positivity: $T$ is completely positive iff $\tau \ge 0$.
        \item Doubly-stochastic: $T(\one)\propto\one$ and $T^*(\one)\propto\one$ iff
              $\tr_A[\tau]\propto \one$ and $\tr_B[\tau]\propto \one$.
        \item Unitality: $T(\one)=\one$ iff $\tr_B[\tau]=\one_{d'}/d$.
        \item Preservation of the trace: $\tr_A[\tau] = (T^*(\one))^{T}/d$,
              i.e., $T^*(\one)=\one$ iff $\tr_A[\tau]=\one_d/d$.
        \item Normalization: $\tr[\tau]=\tr[T^*(\one)]/d$.
    \end{itemize}
\end{proposition}

\begin{proof}
    The complete positivity correspondence was proven in Prop.~1.2 and the other
    equivalences follow from \eqref{eq:2.1} by direct inspection. So, the only piece to
    prove is that the two relations in \eqref{eq:2.1} are mutual inverses which then
    establishes the claimed one-to-one correspondence. To this end let us start with the
    rightmost expression in \eqref{eq:2.1} and insert $\tau=(T\otimes \id)(\ketbra{\Omega}{\Omega})$:
    \begin{align}
        d\,\tr\![\tau\,(A \otimes B^{T})]
         &= \tr\![F^{T_B}\,T^*(A)\otimes B^{T}], \label{eq:2.2}\\
         &= \tr\![F\,T^*(A)\otimes B]
        = \tr\![A\,T(B)], \label{eq:2.3}
    \end{align}
    where we have used the basic tools from Example~1.2, in particular that
    $d\ketbra{\Omega}{\Omega}=F^{T_B}$. Eqs.~\eqref{eq:2.2}--\eqref{eq:2.3} show that
    indeed $T \to \tau \to T$. So the two maps in \eqref{eq:2.1} are mutual inverses if
    $T \mapsto \tau$ is surjective. This is, however, easily seen by decomposing $\tau$
    into a linear combination of rank-one operators $\ket{\psi}\!\bra{\psi'}$ and using
    $\ket{\psi}=(X\otimes \one)\ket{\Omega}$ from Eq.~(1.13).
\end{proof}

When writing out the mapping $\tau \to T$ from \eqref{eq:2.1} in terms of a product
basis we get
\begin{equation}\label{eq:2.4}
    T(B) = d \sum_{i j k l} \bra{ij}\tau\ket{kl}  \ketbra{i}{j}\,B\,\ketbra{l}{k}.
\end{equation}

If $T$ is a quantum channel in the Schr\"odinger picture, then $\tau$ is a density
operator with reduced density matrix $\tr_{B}[\tau] = \one/d$. This means that the
set of quantum channels corresponds one-to-one to the set of bipartite quantum
states which have one reduced density matrix maximally mixed. By replacing $\Omega$
in the construction of $\tau$ by any other pure state we can easily establish a
similar state-channel-duality with respect to states with different reduced density
matrices. Such a correspondence will be one-to-one iff the respective reduced state
has full rank.

If $d$ is countable infinite then $\Omega$ loses its meaning. However, we can partly
restore the above correspondence by using either its unnormalized counterpart
$\sum_{ij}\ketbra{ii}{jj}$ or any non-maximally entangled state which has reduced
density matrix with full rank instead.

The correspondence between $T$ and $\tau$ enables us to show that every linear map
admits a decomposition into at most four completely positive maps:

\begin{proposition}[Decomposition into completely positive maps]
    Every linear map $T \in \mcB(\mcM_d,\mcM_{d'})$ can be written as a complex linear
    combination of four completely positive maps. If $T$ is Hermitian (i.e.,
    $T(B^\dagger)=T(B)^\dagger$ for all $B$), it can be written as a real linear
    combination of two of them.
\end{proposition}

\begin{proof}
    We use that the mapping $T \leftrightarrow \tau$ in Prop.~2.1 is one-to-one and
    linear so that it suffices to decompose $\tau$. On this level the existence of such
    a decomposition is rather obvious: we can always decompose $\tau$ into a Hermitian
    and anti-Hermitian part and each of them into a positive and negative part. More
    explicitly,
    \[
        \tau = \frac{\tau+\tau^\dagger}{2} + i\,\frac{i\tau^\dagger - i\tau}{2}
    \]
    provides a complex linear combination of two Hermitian parts and each of the latter
    can, by invoking the spectral decomposition, be written as a real linear combination
    of two positive semi-definite matrices.
\end{proof}

As a first simple application of this (extremely useful and later extensively used)
correspondence we show the following:

\begin{proposition}[No information without disturbance]
    Consider an instrument represented by a set of completely positive maps
    $\{T_\alpha:\mcM_d \to \mcM_d\}$. If there is no disturbance on average, i.e.,
    $T=\sum_\alpha T_\alpha$ satisfies $T=\id$, then $T_\alpha \propto \id$ for every
    $\alpha$ and the probability for obtaining an outcome $\alpha$ (given by
    $\tr[T_\alpha(\rho)]$) is independent of the input $\rho$ (hence, no information gain).
\end{proposition}

\begin{proof}
    On the level of the Jamiolkowski states the decomposition $\id=\sum_\alpha T_\alpha$
    reads $\ketbra{\Omega}{\Omega}=\sum_\alpha \tau_\alpha$. Since $\tau_\alpha \ge 0$
    (due to complete positivity of $T_\alpha$) this corresponds to a convex decomposition
    of the density operator for the maximally entangled state. Since the latter is pure
    there is only the trivial decomposition, so that $\tau_\alpha = c_\alpha\ketbra{\Omega}{\Omega}$
    for some constants $c_\alpha \ge 0$. Consequently, $T_\alpha = c_\alpha\id$ so
    that $\tr[T_\alpha(\rho)] = c_\alpha$ independent of $\rho$.
\end{proof}

\paragraph{Implementation via teleportation.}
If $T$ is a quantum channel then its Jamiolkowski state $\tau$ can operationally be
obtained by letting the channel act on a maximally entangled state. What about the
converse? If $\tau$ is given, does this help us to implement $T$ as an action on an
arbitrary input $\rho$?

The answer to this involves some form of teleportation: assume that the bipartite
state $\tau$ is shared by two parties, Alice and Bob, so that Bob has the maximally
mixed reduced state. Suppose that Bob has an additional state $\rho$ and that he
performs a measurement on his composite system using a POVM which contains the
maximally entangled state $\omega := \ketbra{\Omega}{\Omega}$ as an effect operator.
We claim now that Alice's state is given by $T(\rho)$ whenever Bob has obtained a
measurement outcome corresponding to $\omega$. In order to show this denote Alice's
reduced density matrix (in case of Bob's success) by $\rho_A$, the success
probability by $p$ and compute the expectation value with an arbitrary operator $A$:
\begin{align}
    p\,\tr[A\rho_A]
     &= \tr[(\tau \otimes \rho)(A \otimes \omega)], \label{eq:2.5}\\
     &= \frac{1}{d} \tr[(\tau \otimes \rho^{T})(A \otimes F)], \label{eq:2.6}\\
     &= \frac{1}{d} \tr[\tau\,(A \otimes \rho^{T})]
    = \frac{1}{d^2} \tr[A\,T(\rho)], \label{eq:2.7}
\end{align}
where we used the basic ingredients from Example~1.2 again and, in the last step,
the r.h.s.\ of \eqref{eq:2.1}. This shows that the described protocol is indeed
successful with probability $p=1/d^2$. In some cases Alice and Bob can, however, do
much better: assume that there is a set of local unitaries
$\{V_i\otimes U_i\}_{i=1}^N$ w.r.t.\ which
$\tau=(V_i\otimes U_i)\tau(V_i\otimes U_i)^\dagger$ is invariant and which are
orthogonal in the sense that $\tr[U_i U_j^\dagger]=d\,\delta_{ij}$. Due to the latter
condition Bob can use a POVM which contains all
$(\one \otimes U_i)^\dagger \omega (\one\otimes U_i)$ as effect operators. When Bob
obtains one of the corresponding outcomes $i$ and he reports this to Alice, she can
``undo'' the action of the unitary by applying $V_i$ with the consequence that
$\rho_A = T(\rho)$ with probability $p=N/d^2$. So if the $U_i$'s form a basis (i.e.,
$N=d^2$) the two protagonists can implement $T$ with unit probability by
``teleporting'' through $\tau$. If $T=\id$ and therefore $\tau=\ketbra{\Omega}{\Omega}$
this reduces to the standard protocol for entanglement-assisted teleportation.

Note that by linearity the implemented action of $T$ does not only hold for
uncorrelated states but it also works properly if $\rho$ is part of a larger system.

\setcounter{problem}{5}
\begin{problem}[Implementation via teleportation]
Suppose Alice and Bob are given the Jamiolkowski state $\tau$ corresponding to a
given quantum channel $T$. Which is the largest probability at which Bob can
teleport an unknown state $\rho$ to Alice (using $\tau$, local operations and
classical communication as resource) so that Alice ends up with the state $T(\rho)$?
\end{problem}

\subsection{Kraus, Stinespring and Neumark}

The Choi-Jamiolkowski state-channel duality allows us to translate between
properties of bipartite states and quantum channels. One immediate implication is
a more specific and very useful representation of quantum channels which
corresponds, on the level of the Jamiolkowski state, to a convex (or spectral)
decomposition into rank-one operators:

\begin{theorem}[Kraus representation]
    A linear map $T \in \mcB(\mcM_d,\mcM_{d'})$ is completely positive iff it admits a
    representation of the form
    \begin{equation}\label{eq:2.8}
        T(A) = \sum_{j=1}^r K_j A K_j^\dagger.
    \end{equation}
    This decomposition has the following properties:
    \begin{enumerate}
        \item Normalization: $T$ is trace preserving iff $\sum_j K_j^\dagger K_j = \one$
              and unital iff $\sum_j K_j K_j^\dagger = \one$.
        \item Kraus rank:\footnote{We call $r=\rank(\tau)$ Kraus rank or Choi rank in order
                  not to confuse it with the rank of $T$ as a linear map. Take for instance $T=\id$
                  the ideal channel. As this is obviously invertible it has full rank as a linear
                  map. However, its Kraus rank is $r=1$. $T(B)=\tr[B]$ instead has rank one but
                  Kraus rank equal to $d$, the dimension of the input space.}
              The minimal number of Kraus operators
              $\{K_j \in \mcB(\CC^d,\CC^{d'})\}_{j=1}^r$ is
              $r=\rank(\tau)\le dd'$.
        \item Orthogonality: There is always a representation with $r=\rank(\tau)$
              Hilbert--Schmidt orthogonal Kraus operators (i.e., $\tr[K_i^\dagger K_j]\propto \delta_{ij}$).
        \item Freedom: Two sets of Kraus operators $\{K_j\}$ and $\{\widetilde{K}_\ell\}$
              represent the same map $T$ iff there is a unitary $U$ so that
              $K_j=\sum_\ell U_{j\ell}\widetilde{K}_\ell$ (where the smaller set is padded with zeros).
    \end{enumerate}
\end{theorem}

\begin{proof}
    Assume $T$ is completely positive. By Prop.~2.1 this is equivalent to saying that
    $\tau \ge 0$ which allows for a decomposition of the form
    \begin{equation}\label{eq:2.9}
        \tau = \sum_{j=1}^r \ketbra{\psi_j}{\psi_j}
        = \sum_{j=1}^r (K_j \otimes \one) \ketbra{\Omega}{\Omega}\,(K_j \otimes \one)^\dagger,
    \end{equation}
    where the first step uses $\tau \ge 0$ and the second Eq.~(1.13). Comparing the
    r.h.s.\ of \eqref{eq:2.9} with the definition $\tau := (T \otimes \id)(\ketbra{\Omega}{\Omega})$
    and recalling that $T \leftrightarrow \tau$ is one-to-one leads to the desired
    decomposition in \eqref{eq:2.8}. It also shows that $r \ge \rank(\tau)$ where
    equality can be achieved and if the $\psi_j$'s are in addition chosen to be
    orthogonal, then the Kraus operators are orthogonal (w.r.t.\ the Hilbert--Schmidt
    inner product) as well.

    Conversely, if a map is of the form \eqref{eq:2.8} then $\tau \ge 0$ which implies
    complete positivity. The conditions for unitality and preservation of the trace are
    straight forward. It remains to show that the freedom in the representation is
    precisely given by unitary linear combinations. This is a direct consequence of
    \eqref{eq:2.9} and the subsequent proposition.
\end{proof}

\begin{proposition}[Equivalence of ensembles]
    Two ensembles of (not necessarily normalized) vectors $\{\psi_j\}$ and
    $\{\widetilde{\psi}_\ell\}$ satisfy
    \begin{equation}\label{eq:2.10}
        \sum_j \ketbra{\psi_j}{\psi_j} = \sum_\ell \ketbra{\widetilde{\psi}_\ell}{\widetilde{\psi}_\ell}
    \end{equation}
    iff there is a unitary $U$ such that
    $\ket{\psi_j} = \sum_\ell U_{j\ell}\ket{\widetilde{\psi}_\ell}$
    (where the smaller set is padded with zero vectors).
\end{proposition}

\begin{proof}
    W.l.o.g.\ we may think of the mixture in \eqref{eq:2.10} as a given density matrix
    $\rho$. From both ensembles we can construct a purification of
    $\rho=\tr_B[\ketbra{\Psi}{\Psi}]$, of the form
    $\ket{\Psi}=\sum_j \ket{\psi_j}\otimes \ket{j}$ where $\{\ket{j}\}$ is an orthonormal
    basis which we use as well in
    $\ket{\widetilde{\Psi}}=\sum_\ell \ket{\widetilde{\psi}_\ell}\otimes \ket{\ell}$.
    It follows from the Schmidt decomposition (Prop.~1.1) that two purifications differ
    by a unitary (or isometry if the dimensions do not match), i.e.,
    $\ket{\Psi}=(\one\otimes U)\ket{\widetilde{\Psi}}$. By taking the scalar product
    with a basis vector $\bra{j}$ on the second tensor factor this leads to
    $\ket{\psi_j}=\sum_\ell U_{j\ell}\ket{\widetilde{\psi}_\ell}$ which proves necessity
    of the condition.\footnote{If the $U$ relating the two purifications is an isometry
        it can always be embedded into a unitary, just by completing the set of orthonormal
        row or column vectors to an orthonormal basis.} Sufficiency follows by inspection
    from unitarity of $U$.
\end{proof}

Note that the number of linearly independent Kraus operators is $r=\rank(\tau)$
independent of the representation. The Kraus representations of a completely
positive map $T$ and its dual $T^*$ are related via interchanging
$K_j \leftrightarrow K_j^\dagger$. Moreover, $T^* = T$ iff there is a
representation with Hermitian Kraus operators (a simple exercise using the unitary
freedom).

\paragraph{Stinespring and the open-system point of view.}
A common perspective is to regard a quantum channel as a description of
``open system dynamics'' which originates from considering only parts of a unitarily
evolving system. Using the Kraus representation we will now see that indeed every
quantum channel can be represented as arising in this way. Before we discuss this
in the Schr\"odinger picture (which might be the more familiar version for
physicists), we will state the equivalent theorem in the Heisenberg picture (more
common to operator algebraists):

\begin{theorem}[Stinespring representation]
    Let $T:\mcM_d \to \mcM_{d'}$ be a completely positive linear map. Then for every
    $r \ge \rank(\tau)$ (recall that $\rank(\tau)\le dd'$) there is a
    $V:\CC^d \to \CC^{d'} \otimes \CC^r$ such that
    \begin{equation}\label{eq:2.11}
        T^*(A) = V^\dagger (A \otimes \one_r)V, \qquad \forall A \in \mcM_{d'}.
    \end{equation}
    $V$ is an isometry (i.e., $V^\dagger V = \one_d$) iff $T$ is trace preserving.
\end{theorem}

\begin{proof}
    This is a simple consequence of the Kraus representation
    $T^*(A)=\sum_{j=1}^r K_j^\dagger A K_j$: construct
    $V := \sum_j K_j \otimes \ket{j}$ with $\{\ket{j}\}$ any orthonormal basis in
    $\CC^r$. Then \eqref{eq:2.11} holds by construction, $r$ is at least the Kraus rank
    (hence minimally $r=\rank(\tau)$, see Thm.~2.1) and
    $V^\dagger V=\sum_j K_j^\dagger K_j=\one$ just reflects the trace preserving
    condition.
\end{proof}

Note from the proof that a representation of the form \eqref{eq:2.11} exists for $T$
as well (and not only for $T^*$). However, the trace-preserving condition is easier
expressed in terms of $T^*$ where it becomes unitality.

The ancillary space $\CC^r$ is usually called \emph{dilation space}. In the same way
in which we constructed $V$ from the set of Kraus operators, we can obtain the
latter from $V$ as $K_j = (\one_{d'}\otimes \bra{j})V$. As $r=\rank(\tau)$ is the
smallest number of Kraus operators, it is also the least possible dimension for a
representation of the form \eqref{eq:2.11}. Dilations with $r=\rank(\tau)$ are
called \emph{minimal}. From the unitary freedom in the choice of Kraus operators
(Thm.~2.1) we obtain that for minimal dilations $V$ is unique up to the obvious
unitary freedom $V \to (\one_{d'} \otimes U)V$. All other dilations can be obtained
by an isometry $U$.

An alternative, but equivalent, way of characterizing minimal dilations is the
identity
\begin{equation}\label{eq:2.12}
    \CC^{d'} \otimes \CC^r
    = \{ (A \otimes \one_r)V\ket{\psi} \mid A\in \mcM_{d'},\ \psi\in\CC^d \},
\end{equation}
i.e., the requirement that the set on the r.h.s.\ spans the entire space.

There is a natural partial order in the set of completely positive maps: we write
$T_2 \ge T_1$ iff $T_2-T_1$ is completely positive. Due to the Choi-Jamiolkowski
representation (Prop.~2.1) this is equivalent to $\tau_2 \ge \tau_1$. Such an order
relation is reflected in the Stinespring representation in the following way:

\begin{theorem}[Relation between ordered cp-maps]
    Let $T_i:\mcM_{d'} \to \mcM_d$, $i=1,2$ be two completely positive linear maps with
    $T_2 \ge T_1$. If $V_i:\CC^d \to \CC^{d'} \otimes \CC^{r_i}$ provide Stinespring
    representations $T_i(A)=V_i^\dagger(A\otimes \one_{r_i})V_i$, then there is a
    contraction $C:\CC^{r_2} \to \CC^{r_1}$ such that $V_1=(\one_{d'}\otimes C)V_2$. If
    $V_2$ belongs to a minimal dilation, then $C$ is unique (for given $V_1$ and $V_2$).
\end{theorem}

\begin{proof}
    We exploit the fact that $T_2 \ge T_1$ is equivalent to $\tau_2 \ge \tau_1$, i.e.,
    the analogous order relation between the corresponding Choi-Jamiolkowski operators.
    Define
    \[
        W_i := (\one_{r_i} \otimes \bra{\Omega}) (V_i \otimes \one_{d'})
        \in \mcB(\CC^d\otimes \CC^{d'},\CC^{r_i}).
    \]
    Then for all $\psi \in \CC^d \otimes \CC^{d'}$,
    \begin{equation}\label{eq:2.13}
        \|W_2\ket{\psi}\|^2 = \bra{\psi}\tau_2\ket{\psi} \ge \bra{\psi}\tau_1\ket{\psi}
        = \|W_1\ket{\psi}\|^2.
    \end{equation}
    Hence, there is a contraction (meaning $C^\dagger C \le \one$)
    $C:\CC^{r_2}\to \CC^{r_1}$ such that $W_1 = C W_2$. Since the map $V_i\mapsto W_i$
    is one-to-one\footnote{The inverse is
        $V_i = d'^2 (W_i\otimes \one_{d'})(\one_d\otimes \ket{\Omega})$.}
    this implies $V_1=(\one_{d'}\otimes C)V_2$. Moreover, if $r_2=\rank(\tau_2)$ (i.e.,
    $\CC^{r_2}$ is a minimal dilation space), then $W_2$ must be surjective so that $C$
    becomes uniquely defined.
\end{proof}

As a simple corollary of this result one obtains a Radon--Nikodym type theorem for
instruments:

\begin{theorem}[Radon--Nikodym for quantum instruments]
    Let $\{T_i\}$ be a set of completely positive linear maps such that
    $\sum_i T_i = T \in \mcB(\mcM_{d'},\mcM_d)$ with Stinespring representation
    $T(A)=V^\dagger(A\otimes \one_r)V$. Then there is a set of positive operators
    $P_i \in \mcM_r$ which sum up to $\one_r$ such that $T_i(A)=V^\dagger(A\otimes P_i)V$.
\end{theorem}

Now let us turn to the Schr\"odinger picture again and see how to represent a
quantum channel from a system-plus-environment point of view:

\begin{theorem}[Open-system representation]
    Let $T:\mcM_d \to \mcM_{d'}$ be a completely positive and trace-preserving linear
    map. Then there is a unitary $U\in \mcM_{d(d')^2}$ and a normalized vector
    $\varphi \in \CC^{d'} \otimes \CC^{d'}$ such that
    \begin{equation}\label{eq:2.14}
        T(\rho) = \tr_E\!\left[U(\rho \otimes \ketbra{\varphi}{\varphi})U^\dagger\right],
    \end{equation}
    where $\tr_E$ denotes the partial trace over the first two tensor factors of the
    involved Hilbert space $\CC^d \otimes \CC^{d'} \otimes \CC^{d'}$.
\end{theorem}

\begin{proof}
    By expressing Stinespring's representation theorem in the Schr\"odinger picture we
    get $T(\rho)=\tr_{\CC^r}[V\rho V^\dagger]$. Let us choose a dilation space of
    dimension $r=dd'$. In this way we can embed $V$ into a unitary which acts on a
    tensor product and write $V = U(\one_d \otimes \ket{\varphi})$ for some
    $\varphi \in \CC^{d'}\otimes \CC^{d'}$ so that \eqref{eq:2.14} follows.
\end{proof}

Notice the possible departure from the dimension of the minimal dilation space in
Thm.~2.5: if $d$ is not a factor of $d'\rank(\tau)$, then the construction fails and
we have to use a space of larger dimension (e.g., $r=dd'$ which is always possible).
In this case the type of freedom in the representation is less obvious.

The physical interpretation of Thm.~2.5 is clear: we couple a system to an
environment, which is initially in a state $\varphi$, let them evolve jointly
according to a unitary $U$ and finally disregard (i.e., trace out) environmental
degrees of freedom. This way of representing quantum channels nicely reminds us of
the fundamental assumption used when we express an evolution in terms of a linear,
completely positive, and trace preserving map: the initial state of the system has
to be independent of the ``environment''---in other words $T$ itself must not depend
on the input $\rho$.

Let us finally revisit Thm.~2.4 from the system-plus-environment point of view:

\begin{proposition}[Environment induced instruments]
    Let $T:\mcM_d \to \mcM_{d'}$ be a completely positive and trace-preserving linear
    with a system-plus-environment representation as in \eqref{eq:2.14}. For every
    decomposition of the form $T=\sum_i T_i$ into completely positive maps $T_i$ there
    is a POVM $\{P_i \in \mcM_{dd'}\}$ such that
    \begin{equation}\label{eq:2.15}
        T_i(\rho) = \tr_E\!\left[(P_i \otimes \one_{d'})\,U(\rho \otimes \ketbra{\varphi}{\varphi})U^\dagger\right].
    \end{equation}
    The Kraus rank $k_i$ of $T_i$ satisfies $k_i \le \rank(P_i)$.
\end{proposition}

\begin{proof}
    \dots\ wouldn't be necessary, as this can essentially be considered a rewriting of
    Thm.~2.4. However, we want to provide a simple alternative proof which is not based
    on Thm.~2.4. Consider the Jamiolkowski state $\tau$ of $T$ and its purification
    \[
        \ket{\psi}:=(\one_d\otimes U)(\ket{\Omega}\otimes \ket{\varphi})
    \]
    so that $\tau=\tr_E[\ketbra{\psi}{\psi}]$. The decomposition of $T$ corresponds to a
    decomposition $\tau=\sum_i \tau_i$ on the level of Choi-Jamiolkowski operators. As
    this is a convex decomposition of the reduced density matrix of $\psi$ we can invoke
    the quantum steering result in Prop.~1.3 (together with the one-to-one correspondence
    between $T_i$ and $\tau_i$) in order to arrive at \eqref{eq:2.15}. From this the
    bound on the Kraus rank follows by utilizing $k_i=\rank(\tau_i)$.
\end{proof}

\paragraph{Neumark's representation of POVMs.}
In the same way as Stinespring's theorem enables us to regard a quantum channel as
part of larger, unitarily evolving system, one can represent a POVM as a von Neumann
measurement performed in an extended space. To simplify matters let us assume that
the effect operators all have rank one (which can always be achieved by spectral
decomposition).

\begin{theorem}[Neumark's theorem]\footnote{As always we restrict ourselves to the
        case of finite dimensions and finite outcomes. Neumark's theorem holds, however,
        similarly if the set of measurement outcomes is characterized by any regular,
        positive, $\mcB(\mathcal{H})$-valued measure on a compact Hausdorff space. The general
        theorem can be viewed as a corollary of Stinespring's representation theorem.}
    Consider a POVM with effect operators $\{\ketbra{\psi_i}{\psi_i}\}_{i=1}^n$ acting on
    $\CC^d$, i.e., $\sum_i \ketbra{\psi_i}{\psi_i} = \one_d$. There exists an orthonormal
    basis $\{\phi_i\}_{i=1}^n$ in $\CC^n \supseteq \CC^d$ such that each $\psi_i$ is the
    restriction of $\phi_i$ to $\CC^d$.
\end{theorem}

\begin{proof}
    Take an orthonormal basis $\{\ket{j}\in\CC^n\}_{j=1}^n$ so that $\CC^d$ is embedded as
    a subspace spanned by the vectors $\ket{1},\dots,\ket{d}$. Then the matrix
    $\Psi \in M_{n,d}$ with components $\Psi_{i,j}:=\braket{j}{\psi_i}$ satisfies
    $\Psi^\dagger \Psi=\one_d$. That is, $\Psi$ is an isometry and can be extended to a
    unitary in $M_n$ by completing the set of orthonormal column vectors to an
    orthonormal basis. The vectors $\phi_j$ are then obtained from the extended
    $\Psi \in M_n$ by $\braket{j}{\phi_i}:=\Psi_{i,j}$.
\end{proof}

Note that in order to implement a general POVM with effect operators $\{P_i\}_{i=1}^n$
as von Neumann measurement the construction leading to Thm.~2.6 requires a space of
dimension $\sum_i \rank(P_i)$. Sometimes we can, however, do better by using rank-one
projections which appear in the decomposition of more than one $P_i$.

\begin{problem}[Minimal Neumark dilation]
Which is the minimal dimension required to implement a given POVM as von Neumann
measurement in an extended space if post-processing of the measurement outcomes
(such as introducing randomness) is allowed?
\end{problem}

Neumark's theorem shows that one can always implement a POVM (at least in
principle) as a von Neumann/sharp measurement in a larger space. That is, we embed
the system which is, say, described by a density matrix $\rho$ into a larger space
where the additional levels are just not populated so that the state becomes a
direct sum (i.e., block diagonal) $\rho \oplus 0$ and we obtain the identity
\begin{equation}\label{eq:2.16}
    \bra{\psi_i}\rho\ket{\psi_i} = \bra{\phi_i}(\rho \oplus 0)\ket{\phi_i}.
\end{equation}
For actual practical implementations a tensor product would be more convenient than
a direct sum structure, since this would allow a realization by coupling the system
to an ancillary system. Such an ``ancilla representation'' of a POVM can easily be
obtained from Neumark's theorem by further enlarging the space (if necessary). Let
us extend the space until the dimension is a multiple of $d$, say $dd' \ge n$. As the
zero-block now has dimension $(d'-1)d$, this enables us to write
$\rho\oplus 0 = \rho \otimes \ketbra{1}{1}$ (with $\ket{1}$ an element of the
computational basis in $\CC^{d'}$). Embedding the vectors $\phi_i$ into this space,
such that $\ket{\Phi_i}:=\ket{\phi_i}\oplus 0$ becomes part of an orthonormal basis in
$\CC^d\otimes \CC^{d'}$, we get
\begin{equation}\label{eq:2.17}
    \bra{\psi_i}\rho\ket{\psi_i} = \bra{\Phi_i}(\rho\otimes \ketbra{1}{1})\ket{\Phi_i}.
\end{equation}

\subsection{Linear maps as matrices}

The last representation, which we discuss, is the simplest one---linear maps
represented as matrices in a way that concatenation of maps corresponds to
multiplication of the respective matrices.

To this end, we use that the space $M_{d,d'}(\CC)$ of complex valued $d\times d'$
matrices is a vector space, i.e., it is closed under linear combinations and scalar
multiplication. As a vector space $M_{d,d'}(\CC)$ is isomorphic to $\CC^{dd'}$. We
can upgrade it to a Hilbert space by equipping it with a scalar product. A common
choice is
\begin{equation}\label{eq:2.18}
    \langle A,B\rangle := \tr\![P\,A^\dagger B], \qquad A,B \in M_{d,d'},
\end{equation}
where $P\in \mcM_{d'}$ is any positive definite operator. Each set of $dd'$
operators which is orthonormal w.r.t.\ this scalar product
($\langle A_i,A_j\rangle=\delta_{ij}$, $i,j=1,\dots,dd'$) forms a basis of
$M_{d,d'}$. Such orthonormal bases lead to simple completely positive maps of the
form
\begin{equation}\label{eq:2.19}
    \sum_{i=1}^{dd'} A_i^\dagger \rho A_i = \tr[\rho]\,P^{-1}, \qquad \forall \rho \in \mcM_d.
\end{equation}
If $P=\one$ in \eqref{eq:2.18} the scalar product is called \emph{Hilbert--Schmidt
    inner product} and the respective space \emph{Hilbert--Schmidt Hilbert space}. Unless
otherwise stated we will in the following use the Hilbert--Schmidt inner product.

The space of linear maps of the form $T:\mcM_d \to \mcM_{d'}$ is isomorphic to
$M_{(d')^2,d^2}$. A concrete matrix representation, which we denote by
$\widehat{T}\in M_{(d')^2,d^2}$ and refer to as \emph{transfer matrix}, is then
obtained in given orthonormal bases $\{G_\beta \in \mcM_d\}_{\beta=1}^{d^2}$ and
$\{F_\alpha \in \mcM_{d'}\}_{\alpha=1}^{(d')^2}$ via
\begin{equation}\label{eq:2.20}
    \widehat{T}_{\alpha,\beta} := \tr\![F_\alpha^\dagger T(G_\beta)].
\end{equation}
By construction a concatenation of maps, e.g.,
$T''(\rho):=T'(T(\rho))$, then corresponds to a matrix multiplication
$\widehat{T''}=\widehat{T'} \widehat{T}$. With some abuse of notation we usually
write the composed map as a formal product $T'' = T' T$ as well.

From \eqref{eq:2.20} we can see that the matrix representation of $T^*$ is given
by the Hermitian conjugate matrix $\widehat{T}^\dagger$ if $T$ is a Hermitian map
(i.e., $T(X^\dagger)=T(X)^\dagger$) and by the transpose matrix $\widehat{T}^{T}$
if the bases are Hermitian.

\begin{proposition}[Self-dual channels]
    Let $T:\mcM_d \to \mcM_d$ be a completely positive linear map. The following are
    equivalent:
    \begin{enumerate}
        \item $T = T^*$.
        \item $\widehat{T}=\widehat{T}^\dagger$ if identical bases $F_\alpha=G_\alpha$ are chosen.
        \item there exists a set of Hermitian Kraus operators for $T$.
    \end{enumerate}
\end{proposition}

\begin{proof}
    $1\leftrightarrow 2$ follows from direct inspection. Also $3\to 1$ is obvious as the
    Kraus operators of $T$ and $T^*$ differ by Hermitian conjugation. In order to show
    $1\to 3$ we can thus write
    \[
        T(X) = \sum_j K_j X K_j^\dagger
        = \frac{1}{2}\sum_j(K_j X K_j^\dagger + K_j^\dagger X K_j).
    \]
    Every pair $(K_j,K_j^\dagger)/\sqrt{2}$ of Kraus operators in the last representation
    can, however, be transformed into a pair of Hermitian operators
    $(K_j+K_j^\dagger,\ i(K_j-K_j^\dagger))/2$ by a unitary linear combination.
    This provides a different representation of the same map (see Prop.~2.1, item~4.)
    by Hermitian Kraus operators.
\end{proof}

Let us in the following have a closer look at different operator bases. The simplest
basis which is orthonormal w.r.t.\ the Hilbert--Schmidt inner product (i.e.,
$\tr[G_\beta^\dagger G_{\beta'}]=\delta_{\beta,\beta'}$) is given by matrix units of
the form $G_\beta=\ketbra{k}{\ell}$ where $\beta$ is identified with the index pair
$(k,\ell)=:\beta$ and $k,\ell=1,\dots,d$. Making this identification explicit leads
to the isomorphism $\mcM_d \simeq \CC^{d^2}$ via the correspondence
$\ketbra{k}{\ell} \leftrightarrow \ket{k,\ell}$. When applied to an arbitrary linear
map represented as $T(\cdot)=\sum_j L_j\,\cdot\,R_j$ this gives
\begin{equation}\label{eq:2.21}
    \widehat{T} = \sum_j L_j \otimes R_j^{T}.
\end{equation}
That is, in particular for completely positive maps with Kraus representation
$T(\cdot)=\sum_j K_j\,\cdot\,K_j^\dagger$ we obtain the simple expression
$\widehat{T}=\sum_j K_j \otimes \overline{K_j}$, which is independent of the
particular choice of Kraus operators (as the unitaries relating the different sets
of Kraus operators cancel). Another advantage of using matrix units as a basis is a
simple mapping between $\widehat{T}$ and the Choi-Jamiolkowski operator
$\tau=(T\otimes \id)(\ketbra{\Omega}{\Omega})$ for arbitrary linear maps:
\begin{equation}\label{eq:2.22}
    \widehat{T} = d\,\tau^\Gamma,
\end{equation}
where $\tau\mapsto \tau^\Gamma$ is an involution defined by
$\bra{m,n}\tau^\Gamma\ket{k,\ell} := \bra{m,k}\tau\ket{n,\ell}$.

\paragraph{More operator bases.}
Despite the useful relations in Eqs.~\eqref{eq:2.21}--\eqref{eq:2.22} it is
sometimes advantageous to use operator bases other than matrix units. The most
common ones can be regarded as generalization of the $2\times 2$ Pauli matrices
(incl.\ identity matrix) to higher dimensions. Since the two-dimensional case is
rather special one has to drop some of the nice properties---for instance
generalize to Hermitian or unitary bases.

\begin{problem}[Generalizing Pauli matrices]
Construct Hilbert--Schmidt orthogonal bases of operators in $\mcM_d$ (different from
tensor products of Pauli matrices) which are unitary and Hermitian.
\end{problem}

\paragraph{Hermitian operator bases.}
The set of Hermitian matrices forms a real vector space. Thus, an orthonormal basis
of Hermitian operators helps to have this property nicely reflected in a concrete
representation. A simple example for such a basis can be constructed by embedding
normalized Pauli matrices $\sigma_x/\sqrt{2}$ and $\sigma_y/\sqrt{2}$ as $2\times 2$
principal submatrices into $\mcM_d$ (so that only two entries are non-zero). This
leads to $d^2-d$ orthonormal matrices. In order to complete the basis we add $d$
diagonal matrices which we construct from any orthogonal matrix $M\in M_d(\RR)$ by
choosing the diagonal entries of the $k$'th diagonal matrix equal to the $k$'th
column vector of $M$. This yields a complete Hilbert--Schmidt orthonormal basis. If
in addition one column of $M$ leads to $\one_d$, then the other $d^2-1$ matrices are
traceless and generate the Lie-algebra $\mathfrak{su}(d)$, i.e., they provide a
complete set of infinitesimal generators of $\mathrm{SU}(d)$. A popular example of
this kind are the $8$ Gell--Mann matrices in $\mcM_3$ (and, of course, the $3$ Pauli
matrices in $\mcM_2$).

\paragraph{Unitary operator bases.}
Unitary operator bases play an important role in quantum information theory as they
correspond one-to-one to bases of maximally entangled states (see Example~1.2). In
$\mcM_2$ all orthogonal bases of unitaries are of the form
\begin{equation}\label{eq:2.23}
    U_k = e^{i\phi_k} V_1 \sigma_k V_2, \qquad k=0,\dots,3,
\end{equation}
where $V_1$ and $V_2$ are unitaries and $\sigma_k$ denote the Pauli matrices (with
$\sigma_0=\one$). That is, in $\mcM_2$ all orthogonal bases of unitaries are
essentially equivalent to the set of Pauli matrices. In higher dimensions a similar
characterization of all unitary operator bases is not known. A priori, it is quite
remarkable that such orthogonal bases exist in every dimension. Note that this
depends crucially on the chosen scalar product: if we choose the $A_i$'s in
\eqref{eq:2.19} proportional to unitaries (and thus $d=d'$), then $\rho=\one$ shows
that they can only form an orthogonal basis if $P \propto \one$, corresponding to
the Hilbert--Schmidt scalar product.

\begin{example}[Discrete Weyl system and GHZ bases]
    One of the most important unitary operator bases and also a very inspiring example
    regarding the construction of others comes from a discrete version of the
    Heisenberg--Weyl group. Consider a set
    $\{U_{kl}\in \mcM_d\}_{k,\ell=0}^{d-1}$ of $d^2$ unitaries defined by
    \begin{equation}\label{eq:2.24}
        U_{kl} := \sum_{r=0}^{d-1} \eta^{r\ell} \ket{k+r}\!\bra{r},\qquad
        \eta := e^{2\pi i/d},
    \end{equation}
    where addition inside the ket is modulo $d$. This set has the following nice
    properties:
    \begin{itemize}
        \item It forms a basis of operators in $\mcM_d$ which is orthogonal w.r.t.\ the
              Hilbert--Schmidt scalar product, i.e., $\tr[U_{ij}^\dagger U_{k\ell}] = d\,\delta_{ik}\delta_{j\ell}$.
              Since $U_{00}=\one$ we have in particular $\tr[U_{k\ell}]=0$ for all
              $(k,\ell)\ne (0,0)$.
        \item It is a discrete Weyl system since $U_{ij}U_{k\ell}=\eta^{jk}U_{i+k,j+\ell}$
              (with addition modulo $d$). Thus $U_{k\ell}^{-1}=\eta^{k\ell}U_{-k,-\ell}$.
        \item The set is generated by $U_{01}$ and $U_{10}$ via $U_{k\ell}=(U_{10})^k(U_{01})^\ell$.
        \item If $d$ is odd, then $U_{k\ell}\in \mathrm{SU}(d)$. For $d$ even
              $\det U_{k\ell} = (-1)^{k+\ell}$.
        \item For $d=2$ the set reduces to the set of Pauli matrices with identity, i.e.,
              $(\sigma_x,\sigma_y,\sigma_z)=(U_{10},\,iU_{11},\,U_{01})$.
        \item The group generated by the unitaries $U_{k\ell}$ is isomorphic to the
              discrete Heisenberg--Weyl group
              \begin{equation}\label{eq:2.25}
                  \left\{\begin{pmatrix}
                      1 &\ell &m\\
                      0 &1    &k\\
                      0 &0    &1
                  \end{pmatrix}
                  \middle|
                  k,\ell,m\in \ZZ_d\right\}.
              \end{equation}
    \end{itemize}

    Among the many applications of this set we will briefly discuss the construction of
    bases of entangled states. As indicated in Example~1.2, $(U_{k\ell}\otimes \one)\ket{\Omega}$
    is an orthonormal basis of maximally entangled states in $\CC^d \otimes \CC^d$. For
    more than two, say $n$, parties a similar construction can be made by exploiting
    some group structure. Consider the group of local unitaries
    \begin{equation}\label{eq:2.26}
        G =
        \left\{\bigotimes_{k=1}^n U_{i,j_k}
        \ \middle|\
        \sum_{k=1}^n j_k \equiv 0 \ (\mathrm{mod}\ d)\right\},
    \end{equation}
    where $i\in \ZZ_d$, $j\in \ZZ_d^n$ and one of the components of $j$, say $j_1$,
    depends on the others by the additional constraint. Utilizing the above properties
    of the set $\{U_{k\ell}\}$ it is readily verified that $G$ is an abelian group with
    $d^n$ elements and that it spans its own commutant $G'$. Since the algebra $G'$ is
    therefore abelian as well, it forms a simplex, i.e., each element in $G'$ has a
    unique convex decomposition into minimal projections. The latter turn out to be
    one-dimensional and can be parametrized by $\omega\in \ZZ_d^n$ such that they
    correspond to vectors of the form
    \begin{equation}\label{eq:2.27}
        \ket{\Psi_\omega} :=
        \frac{1}{\sqrt{d}}\sum_{l=0}^{d-1} \eta^{l\omega_1}\bigotimes_{k=1}^n \ket{\omega_k + l}.
    \end{equation}
    These vectors form an orthonormal basis of the entire Hilbert space
    $(\CC^d)^{\otimes n}$ and they can all be obtained by local unitaries from the GHZ
    state $\ket{\Psi_0}=\sum_k \ket{k}^{\otimes n}/\sqrt{d}$. For $n=d=2$ we get
    $G=\{\one, \sigma_x\otimes \sigma_x,\,-\sigma_y\otimes \sigma_y, \sigma_z\otimes \sigma_z\}$
    (which is the Klein four group\footnote{ok, \emph{``Vierergruppe''} sounds better})
    and the states in \eqref{eq:2.27} become the four Bell states.
\end{example}

Example~2.1 exhibits two properties which suggest possible starting points for the
construction of other unitary operator bases: (i) the group structure of the set,
and (ii) the property that every element is generated by two elementary actions:
shifting the computational basis and multiplying with phases. The latter approach
leads to a fairly general construction scheme. Bases of unitaries obtained in this
way are said to be of shift-and-multiply type, and they are of the general form
\begin{equation}\label{eq:2.28}
    U_{k\ell} = \sum_{r=0}^{d-1} H^{(\ell)}_{kr} \ket{\Lambda_{\ell r}}\!\bra{r},
\end{equation}
where $\{H^{(\ell)}\in \mcM_d\}$ is a set of (not necessarily distinct) complex
Hadamard matrices and $\Lambda \in M_d(\ZZ_d)$ is a Latin square. That is,
$|H^{(\ell)}_{kr}|=1$ are phases so that $H^{(\ell)}(H^{(\ell)})^\dagger=d\one$, and
$\Lambda$ is such that each row and column contains every number from $0$ to $d-1$
exactly once. Obviously, the set in \eqref{eq:2.24} is just the simplest example of a
shift-and-multiply basis. Its group property (the fact, that the elements form a group
of order $d^2$, up to phases) is not shared by all shift-and-multiply bases. For small
dimensions (at least for $d<6$), every unitary basis with such a group structure
(sometimes called \emph{nice error basis}) is of shift-and-multiply type. For higher
dimensions ($d=165$ being the first known counter example) this is no longer valid.

\paragraph{Overcomplete sets, frames and SIC-POVMs.}
We have seen that operator bases for $\mcM_d$ can be constructed from unitary as well
as from Hermitian operators. What about positive semidefinite operators or Hermitian
rank-one projections? This turns out not to be possible if one demands orthogonality
w.r.t.\ the Hilbert--Schmidt scalar product (see Prop.~2.7). However, relaxing this
condition leads to bases with remarkable properties---and our discussion to the
framework of frames. A frame of a vector space, say $\CC^d$, is a set of vectors
$\{\phi_i\}_{i=1}^n$ for which there are constants $0<a\le b<\infty$ such that for all
$\psi\in \CC^d$,
\begin{equation}\label{eq:2.29}
    a\|\psi\|^2 \le \sum_i |\braket{\psi}{\phi_i}|^2 \le b\|\psi\|^2.
\end{equation}
Hence, the $\phi_i$'s have to span the entire vector space. If $a=b$, the frame is
called tight and if in addition $n=d$ it is nothing but an orthonormal basis (with
norm equal to $\sqrt{a}$). Since for tight frames $\sum_i \ketbra{\phi_i}{\phi_i} = a\one$,
they correspond one-to-one to rank-one POVMs.

A POVM is called informationally complete if its measurement outcomes allow for a
complete reconstruction of an arbitrary state. The existence of such POVMs is fairly
obvious: just use properly normalized spectral projections of a unitary or Hermitian
operator basis as effect operators. A simple physical implementation in $\CC^d$
would be to first draw a random number from $1,\dots,d^2$ and then measure the
corresponding element of a Hermitian operator basis. Clearly, a POVM acting on
$\CC^d$ is informationally complete iff it contains $d^2$ linearly independent
effect operators.

Now let us try to construct a positive semidefinite operator basis for $\mcM_d$ which
is ``as orthogonal as possible'' and see how this relates to tight frames and
informationally complete POVMs:

\begin{proposition}[SIC POVMs]
    Let $\{P_i\}_{i=1}^n$ be a set of positive semidefinite operators in $\mcM_d$,
    $d\le n$, which are normalized w.r.t.\ the Hilbert--Schmidt scalar product, i.e.,
    $\langle P_i,P_i\rangle=\tr[P_i^2]=1$. Then
    \begin{equation}\label{eq:2.30}
        \sum_{i\ne j} |\langle P_i,P_j\rangle|^2 \ge \frac{(n-d)^2\,n}{(n-1)d^2},
    \end{equation}
    with equality iff all $P_i$ are rank-one projections fulfilling
    $\langle P_i,P_j\rangle=\frac{n-d}{(n-1)d}$ for all $i\ne j$, and $\sum_i P_i = \frac{n}{d}\one$.
    In case of equality the $P_i$'s are linearly independent.
\end{proposition}

\begin{proof}
    Define $Q:=\sum_i P_i$. Invoking Cauchy--Schwarz inequality twice we get
    \begin{equation}\label{eq:2.31}
        \sum_{i\ne j} |\langle P_i,P_j\rangle|^2 \ge \frac{1}{n^2-n}\left(\sum_{i\ne j} \langle P_i,P_j\rangle\right)^2
    \end{equation}
    \begin{equation}\label{eq:2.32}
        = \frac{1}{n^2-n}(\tr[Q^2]-n)^2 \ge \frac{(\tr[Q]^2/d - n)^2}{n^2-n},
    \end{equation}
    where the first inequality holds iff $\langle P_i,P_j\rangle=\mathrm{const}$ for all
    $i\ne j$ and the second one iff $Q\propto \one$. Since due to normalization
    $\tr[P_i]\ge 1$ and thus $\tr[Q]\ge n$ with equality iff $P_i$ is a rank-one
    projection we can further bound \eqref{eq:2.32} from below leading to the r.h.s.\ of
    \eqref{eq:2.30}. Collecting the conditions for equality completes the proof of the
    first part.

    In order to show linear independence assume $\sum_j c_j P_j=0$ and set
    $c:=\langle P_i,P_j\rangle$ for $i\ne j$. Then for all $i$:
    \[
        0=\sum_j c_j \langle P_i,P_j\rangle = c_i(1-c) + c\sum_j c_j,
    \]
    so $c_i$ is independent of $i$ and therefore $c_i=0$.
\end{proof}

As claimed before, Prop.~2.7 shows that positive semidefinite operators cannot form
a Hilbert--Schmidt orthogonal operator basis (for that the l.h.s.\ of
\eqref{eq:2.30} would have to be zero, which is in conflict with $n=d^2$). Whether
the conditions for equality in Prop.~2.7 are vacuous or actually have a solution
depends on $n$ and $d$. Whenever there exists a solution, it forms a tight frame.
While for $n=d$ the solutions are exactly all orthonormal bases of $\CC^d$, for
$n>d^2$ no solution exists (as the $P_i$'s become necessarily linear dependent then).

Remarkably, there exist solutions for some $n=d^2$, i.e., tight frames which are
informationally complete POVMs with the additional symmetry $\tr[P_iP_j]=(d+1)^{-1}$
for all $i\ne j$. These are called symmetric informationally complete (SIC) POVMs.
For $d=2$ all SIC POVMs are obtained by choosing points on the Bloch sphere which
correspond to vertices of a tetrahedron. For $d=3$ an example can be constructed
from the shift-and-multiply basis in \eqref{eq:2.24} by taking projections onto the
vectors $\ket{\xi_{k\ell}} := U_{k\ell}\ket{\xi_{00}}$ with
$\ket{\xi_{00}} := (\ket{0}+\ket{1})/\sqrt{2}$. In other words, in this example the
vectors onto which the $P_i$'s project are generalized coherent states w.r.t.\ the
discrete Heisenberg--Weyl group. Analogous to continuous coherent states, all SIC
POVMs provide (due to the operator basis property) a representation of arbitrary
operators $\rho\in \mcM_d$:
\begin{equation}\label{eq:2.33}
    \rho = \frac{1}{d}\sum_i\left((d+1)\tr[P_i\rho] - \tr[\rho]\right) P_i.
\end{equation}
In the context of coherent states this is called diagonal representation or (in
quantum optics for the continuous Heisenberg--Weyl group) $P$-function
representation.

Using the effect operators of a SIC-POVM as Kraus operators we get a quantum channel
of the form
\begin{equation}\label{eq:2.34}
    \frac{1}{d}\sum_{i=1}^{d^2} P_i \rho P_i
    = \frac{\one\,\tr[\rho] + \rho}{d+1}, \qquad \forall \rho\in \mcM_d.
\end{equation}
Both \eqref{eq:2.33} and \eqref{eq:2.34} are readily verified by using the
properties (in particular the operator basis property) of the $P_i$'s.

\begin{problem}[SIC POVMs]
Determine the pairs $(n,d)$ for which equality in \eqref{eq:2.30} can be achieved.
\end{problem}

\subsection{Normal forms}

It is sometimes useful to cut some trees in order to get a better overview of the
forest. Depending on the situation will later encounter various ``normal forms'' for
the representations discussed previously---in particular, decompositions based
on spectral, convex or semi-group properties. This section is devoted to normal
forms of completely positive maps w.r.t.\ invertible maps with Kraus rank one.
More specifically, given a completely positive map $T:\mcM_{d_1} \to \mcM_{d_2}$ we
are interested in simple representatives of the equivalence class
\begin{equation}\label{eq:2.35}
    T \sim \Phi_2T\Phi_1,
\end{equation}
where $\Phi_i \in \mcB(\mcM_{d_i})$ is of the form $\Phi_i(\cdot):=X_i\,\cdot\,X_i^\dagger$
with $X_i\in \mathrm{SL}(d_i,\CC)$, i.e., complex valued matrices with unit
determinant. Such maps often run under the name filtering operations. On the level
of Choi-Jamiolkowski operators $\tau\in \mcB(\CC^{d_2}\otimes \CC^{d_1})$ the
transformation in \eqref{eq:2.35} corresponds to
$\tau\mapsto \tau' := (X_2\otimes X_1)\tau (X_2\otimes X_1)^\dagger$
(up to an irrelevant transposition for $X_1$). In order to construct a normal form
w.r.t.\ such transformations we utilize the following optimization problem:
\begin{equation}\label{eq:2.36}
    p := \inf_{X_i\in \mathrm{SL}(d_i,\CC)} \tr[\tau'].
\end{equation}
Recall that we assume complete positivity which by Prop.~2.1 means $\tau\ge 0$.
Clearly, $p\le \tr[\tau]$ (as we can choose $X_i=\one$) and
$p \ge \|X_i\|_\infty^2\,\lambda_{\min}(\tau)$ with $\lambda_{\min}(\tau)$ the
smallest eigenvalue of $\tau$. Our aim is to choose as a representative of the
equivalence class one which minimizes \eqref{eq:2.36}:

\begin{proposition}[Normal form for generic $\tau$]
    Let $\tau \in \mcB(\CC^{d_2}\otimes \CC^{d_1})$ be positive definite. Then there exist
    $X_i\in \mathrm{SL}(d_i,\CC)$ which attain the infimum in \eqref{eq:2.36} so that the
    respective optimal $\tau'=(X_2\otimes X_1)\tau (X_2\otimes X_1)^\dagger$ is such that
    both partial traces are proportional to the identity matrix.
\end{proposition}

\begin{proof}
    From the upper and lower bound on $p$ we know that we can w.l.o.g.\ restrict to
    $\|X_i\|_\infty^2 \le \tr[\tau]/\lambda_{\min}(\tau)$. These $X_i$ form a compact set
    (since $\tau>0$ by assumption) so that the infimum is indeed attained.

    Eq.~\eqref{eq:2.36} can now be minimized by a simple iteration in $X_1$ and $X_2$:
    denote by $\tau_1=\tr_2[\tau]$ a partial trace of $\tau$ so that
    $\tau'_1 := X_1 \tau_1 X_1^\dagger$ is the corresponding one for
    $\tau':=(\one\otimes X_1)\tau(\one\otimes X_1)^\dagger$. Choosing
    $X_1 := \det(\tau_1)^{1/(2d_1)} \tau_1^{-1/2}$ will lead to $\tau'_1\propto \one$.
    Moreover, the arithmetic-geometric mean inequality gives
    \begin{equation}\label{eq:2.37}
        \tr[\tau'] = \det(\tau_1)^{1/d_1}\,d_1 \le \tr[\tau_1] = \tr[\tau],
    \end{equation}
    with equality iff $\tau_1\propto \one$. Iterating this step w.r.t.\ $X_1$ and $X_2$
    will thus decrease the trace while both reduced density matrices converge to
    something proportional to the identity. As this holds true in particular for the
    optimal $\tau'$, it has to have $\tau'_i \propto \one$.
\end{proof}

Together with the Choi-Jamiolkowski correspondence in Prop.~2.1 this has an
immediate corollary on the level of completely positive maps:

\begin{proposition}[Normal form for generic cp maps]
    Let $T:\mcM_{d_1} \to \mcM_{d_2}$ be a completely positive map with full Kraus rank.
    There exist invertible completely positive maps $\Phi_i \in \mcB(\mcM_{d_i})$ with
    Kraus rank one such that $T':=\Phi_2T\Phi_1$ is doubly-stochastic, i.e., both
    $T(\one)$ and $T^*(\one)$ are proportional to the identity matrix.
\end{proposition}

The normal form in Prop.~2.8 is unique\footnote{This is an exercise about
non-negative matrices: assume $\widetilde{\tau}:=(D_2\otimes D_1)\tau'(D_2\otimes D_1)$
is a second normal form with $D_i$ positive diagonal matrices (we can always use the
unitary freedom to choose them like this). Then $M_{ij}:=\bra{ij}\tau'\ket{ij}$
(and analogously $\widetilde{M}$) is a doubly stochastic non-negative matrix, i.e.,
$M_{ij}\ge 0$ and all rows as well as all columns have the same sum. Moreover,
$\widetilde{M}=M * H$ is a Hadamard product with $H_{ij}:=[D_2]_{ii}[D_1]_{jj}$. The
                only $H$ which achieves such a mapping is, however, a multiple of $H_{ij}=1$.} up to
    local unitaries $\tau' \mapsto (U_2\otimes U_1)\tau'(U_2\otimes U_1)^\dagger$, which
    corresponds to unitary channels $\Phi_i(\cdot)=U_i\,\cdot\,U_i^\dagger$ in
    Prop.~2.9. This implies that the algorithmic procedure used in the proof of
    Prop.~2.8 in order to obtain the normal form by minimizing $\tr[\tau']$ eventually
    converges to the global optimum.

    In cases where $\tau$ has a kernel the infimum in \eqref{eq:2.36} may either be zero
    or not be attained for finite $X_i$. An exhaustive investigation of general normal
    forms w.r.t.\ the equivalence class \eqref{eq:2.35} has been performed for qubit
    channels.

    \paragraph{Qubit maps.}
    We continue the discussion about normal forms w.r.t.\ Kraus-rank one operations for
    completely positive maps $T:\mcM_2 \to \mcM_2$. Choosing normalized Pauli matrices as
    operator basis (with $\sigma_0=\one$) we can represent $T$ as a $4\times 4$ matrix
$\widehat{T}_{ij} := \tr[\sigma_iT(\sigma_j)]/2$. If $T$ is Hermitian, then
$\widehat{T}$ is real, and if $T$ is trace preserving, then
$(\widehat{T}_{0,j})=(1,0,0,0)$. A quantum channel in the Schr\"odinger picture is
    thus represented by
    \begin{equation}\label{eq:2.38}
        \widehat{T} =
        \begin{pmatrix}
            1 &0\\
            v &\Delta
        \end{pmatrix},
    \end{equation}
    where $\Delta$ is a real $3\times 3$ matrix and $v\in \RR^3$. The corresponding
    Jamiolkowski state is given by
    \[
        \tau = \frac{1}{4}\sum_{ij} \widehat{T}_{ij} \sigma_i\otimes \sigma_j^{T},
    \]
    so that $v$ describes its reduced density operator and $\Delta$ its correlations.
    The conditions for $\widehat{T}$ to correspond to a completely positive map have to
    be read off from $\tau\ge 0$. The complexity of this condition is the drawback of
    this representation. One of its advantages is a nice geometric interpretation:
    parameterizing a density operator via $\rho=(\one+\sum_{k=1}^3 x_k\sigma_k)/2$, i.e.,
    in terms of a vector $x\in \RR^3$ within the Bloch ball $\|x\|\le 1$ (see
    Example~1.1), the action of $T$ is a simple affine transformation
    \begin{equation}\label{eq:2.39}
        x \mapsto v + \Delta x.
    \end{equation}
    From here conditions for $T$ to be positive are readily derived (as the vector has
    to stay within the Bloch ball), and we will discuss them in greater detail in
    \dots. The following is a useful proposition in this context:

    \begin{proposition}
        Let $T:M_2(\CC)\to M_2(\CC)$ be a Hermiticity preserving linear map with
        $\Delta_{ij}:=\tr[\sigma_iT(\sigma_j)]/2$ the lower-right submatrix (i.e.,
        $i,j=1,2,3$) of the matrix representation $\widehat{T}$. Then
        $\widehat{T}' := \widehat{T}_{00} \oplus \Delta$ represents a (completely) positive
        map iff $\widehat{T}$ does so.
    \end{proposition}

    \begin{proof}
        We use that $D:=\mathrm{diag}(1,-1,-1,-1)$ is the matrix representation of
        time-reversal (see Sec.~3.3), i.e., matrix transposition in some basis. Following
        the relation in Eq.~(1.24) the map $D\widehat{T}D$ is (completely) positive if
        $\widehat{T}$ is. The same holds thus for the convex combination
        $(\widehat{T}+D\widehat{T}D)/2 = \widehat{T}'$.
    \end{proof}

    Suppose now that we act with a unitary before and another one after applying the
    channel $T$ so that the overall action is
$\rho \mapsto U_2T(U_1\rho U_1^\dagger)U_2^\dagger$.
    The transfer matrix corresponding to this concatenation is then given by the product
$(\one\oplus O_2)\widehat{T}(\one\oplus O_1)$, where $O_i\in \mathrm{SO}(3)$ are real
    rotations of the Bloch sphere. This reflects the two-to-one group homomorphism
$\mathrm{SU}(2)\to \mathrm{SO}(3)$ (see Examples~1.1 and 2.2). We may use this in
    order to diagonalize $\Delta \to \mathrm{diag}(\lambda_1,\lambda_2,\lambda_3)$ so that
$\lambda_1\ge \lambda_2\ge |\lambda_3|$ and for positive, trace-preserving maps
    necessarily $1\ge \lambda_1$. Up to a possibly remaining sign this diagonalization is
    nothing but a real singular value decomposition. Expressed in terms of the $\lambda_i$'s
    a necessary condition for complete positivity is that
    \begin{equation}\label{eq:2.40}
        \lambda_1 + \lambda_2 \le 1 + \lambda_3.
    \end{equation}
    This becomes sufficient if the map is unital, i.e., if $v=0$ (see also Exp.~2.3).

    Extending the freedom in the transformations from $\mathrm{SU}(2)$ to
$\mathrm{SL}(2)$ enables us to further simplify $\widehat{T}$ and to bring it to a
normal form (as stated in Prop.~2.9 for the full rank case):

\begin{proposition}[Lorentz normal form]
    For every qubit channel $T$ there exist invertible completely positive maps
    $\Phi_1,\Phi_2$, both of Kraus rank one, such that the concatenation
    $\Phi_2T\Phi_1 =: T'$ (characterized by $v$ and $\Delta$) is a qubit channel of
    one of the following three forms:
    \begin{enumerate}
        \item Diagonal: $T'$ is unital ($v=0$) with $\Delta$ diagonal. This is the generic
              case (proved in Prop.~2.9).
        \item Non-diagonal: $T'$ has $\Delta=\mathrm{diag}(x/\sqrt{3},\,x/\sqrt{3},\,1/3)$,
              $0\le x \le 1$ and $v=(0,0,2/3)$. These channels have Kraus rank $3$ for $x<1$
              and Kraus rank $2$ for $x=1$.
        \item Singular: $T'$ has $\Delta=0$ and $v=(0,0,1)$. This channel has Kraus rank
              $2$ and is singular in the sense that it maps everything onto the same output.
    \end{enumerate}
\end{proposition}

\begin{example}[Lorentz group and spinor representation]
    Our aim is to understand the action of $T \mapsto \Phi_2T\Phi_1$ as an equivalence
    transformation on $\widehat{T}$ by elements from the Lorentz group. In order to see
    how this arises, consider the space $M_2^\dagger(\CC)$ of complex Hermitian $2\times 2$
    matrices. Every such matrix can be expanded as $M=\sum_{i=0}^3 x_i\sigma_i$ with
    $x\in \RR^4$. Since $\det(M)=x_0^2-x_1^2-x_2^2-x_3^2=\langle x,\eta x\rangle$ with
    $\eta:=\mathrm{diag}(1,-1,-1,-1)$ we can identify $M_2^\dagger(\CC)$ with Minkowski
    space such that the determinant provides the Minkowski metric. If $X\in \mathrm{SL}(2,\CC)$,
    the map
    \begin{equation}\label{eq:2.41}
        M \mapsto XMX^\dagger
    \end{equation}
    in this way becomes a linear isometry in Minkowski space, i.e., a Lorentz
    transformation. In fact, the group $\mathrm{SL}(2,\CC)$ is a double cover of the
    special orthochronous Lorentz group
    \begin{equation}\label{eq:2.42}
        \mathrm{SO}^+(1,3)
        := \{L\in M_4(\RR)\mid \det(L)=1,\ L\eta L^{T}=\eta,\ L_{00}>0\}
    \end{equation}
    in very much the same way as $\mathrm{SU}(2)$ is a double cover of $\mathrm{SO}(3)$.
    The map $\mathrm{SL}(2,\CC)\to \mathrm{SO}^+(1,3)$ constructed above is sometimes
    called spinor map. It is two-to-one since $\pm X$ have the same effect in
    \eqref{eq:2.41}. Due to this equivalence the transfer matrix of the channel
    $\Phi_2T\Phi_1$ becomes
    \begin{equation}\label{eq:2.43}
        L_2\,\widehat{T}\,L_1,\qquad L_i\in \mathrm{SO}^+(1,3).
    \end{equation}
    The normal form in Prop.~2.11 can thus be regarded as a normal form w.r.t.\ special
    orthochronous Lorentz transformations.

    Every $L\in \mathrm{SO}^+(1,3)$ can be decomposed into a ``boost'' and a spatial
    rotation. This decomposition can be obtained from the polar decomposition $X=PU$ in
    $\mathrm{SL}(2,\CC)$, where $P>0$ and $UU^\dagger=\one$. In order to make this more
    explicit, define generators of spatial rotations by
    $R_i := \sum_{j,k=1}^3 \varepsilon_{ijk} \ketbra{k}{j}\in M_4$, and generators of
    boosts by $B_i := \ketbra{0}{i}+\ketbra{i}{0}\in M_4$ with $i=1,2,3$. Then with
    $\vec{n},\vec{m}\in \RR^3$ the mapping $\mathrm{SL}(2,\CC)\to \mathrm{SO}^+(1,3)$
    takes the form
    \begin{equation}\label{eq:2.44}
        U = e^{-i\vec{n}\cdot \vec{\sigma}/2} \mapsto e^{-\vec{n}\cdot \vec{R}},\qquad
        P = e^{\vec{m}\cdot \vec{\sigma}/2} \mapsto e^{\vec{m}\cdot \vec{B}}.
    \end{equation}
    Decomposing $P$ further according to the spectral decomposition leads to a diagonal
    matrix which corresponds to a boost in $z$-direction (i.e., it acts in the
    $0$--$3$--plane as $\one \cosh m_3 + \sigma_x \sinh m_3$).
\end{example}

\begin{example}[Bell diagonal states and Pauli channels]
    \dots
\end{example}

\subsection{Literature}

\end{document}
